\documentclass[10pt]{extarticle}
\usepackage[a4paper,margin=1.5cm]{geometry}
\usepackage[T1]{fontenc}
\usepackage{amsmath,amssymb}
\usepackage{graphicx}
\usepackage[export]{adjustbox}
\usepackage{multicol}
\usepackage{enumitem}
\usepackage{booktabs,array,tabularx}
\setlength{\columnsep}{0.7cm}
\setlength{\parindent}{0pt}
\setlength{\parskip}{3pt}
\pagestyle{empty}
\makeatletter
\newcommand{\qfigdef}[4]{\@namedef{qf@file@#1}{#2}\@namedef{qf@w@#1}{#3}\@namedef{qf@wide@#1}{#4}}
\newcommand{\qf@check}[1]{\@ifundefined{qf@file@#1}{\PackageError{score2}{Unknown figure id #1}{}}{}}
\newcommand{\qf@img}[4][]{\includegraphics[width=#2\dimexpr\@nameuse{qf@w@#4}pt\relax,max width=#3,max height=0.45\textheight,#1]{\@nameuse{qf@file@#4}}}
\newif\ifqf@full
\newif\ifqf@boxed
\newsavebox{\qf@box}
\newcommand{\qf@open}{\global\qf@fulltrue\par\global\qf@boxedtrue
  \setbox\qf@box\vbox\bgroup\hsize\textwidth\linewidth\textwidth}
\newcommand{\qf@emit}[1]{\par\egroup\global\qf@boxedfalse
  \ifdim\dimexpr\ht\qf@box+\dp\qf@box\relax<#1\textheight\noindent\usebox{\qf@box}\par\else\unvbox\qf@box\fi}
\newcommand{\qfullsplit}{\ifqf@boxed\qf@emit{0.85}\fi}
\newcommand{\qf@close}{\ifqf@boxed\qf@emit{0.45}\fi\global\qf@fullfalse}
\newcommand{\qfullbegin}{\end{multicols}\qf@open}
\newcommand{\qfullend}{\qf@close\begin{multicols}{2}}
\newcommand{\qfullstart}{\qf@open}
\newcommand{\qfulldone}{\qf@close}
\newcommand{\qfig}[2][1]{\qf@check{#2}%
  \ifqf@full\qf@blk{#1}{#2}\else
  \ifnum\@nameuse{qf@wide@#2}=1 \expandafter\@firstoftwo\else\expandafter\@secondoftwo\fi
  {\end{multicols}\begin{center}\qf@img{#1}{\textwidth}{#2}\end{center}\begin{multicols}{2}}%
  {\qf@blk{#1}{#2}}\fi}
\newcommand{\qf@blk}[2]{\par\smallskip{\centering\qf@img{#1}{\linewidth}{#2}\par}\smallskip}
\newcommand{\qfigi}[2][1]{\qf@check{#2}\qf@img[valign=c]{#1}{\linewidth}{#2}}
\makeatother
\newcommand{\papersec}[1]{\par\bigskip{\centering\Large\bfseries PAPER-#1\par}\medskip}
\newcommand{\partsec}[2]{\par\medskip\noindent{\large\bfseries PART-#1: #2}\par\smallskip}
\newcommand{\sect}[1]{\par\smallskip\noindent\textbf{#1}\par}
\newcommand{\qn}[1]{\par\smallskip\noindent\textbf{#1.}~}
\newcommand{\paraq}[2]{\par\smallskip\noindent\textbf{Paragraph for Questions #1 and #2}\par}
\newenvironment{qoptions}{\begin{itemize}[leftmargin=1.8em,itemsep=1pt,topsep=2pt,label={}]}{\end{itemize}}
\newcommand{\opt}[2]{\item[(#1)] #2}
% figure registry (generated by score2_tools.py finalize)
\qfigdef{p1_phy_s1_q02_1}{SRG (Enthusiast) - Test 39 (15 Sep 2024)-img/p1_phy_s1_q02_1.png}{173.9}{0}
\qfigdef{p1_phy_s1_q03_1}{SRG (Enthusiast) - Test 39 (15 Sep 2024)-img/p1_phy_s1_q03_1.png}{128.3}{0}
\qfigdef{p1_phy_s1_q04_1}{SRG (Enthusiast) - Test 39 (15 Sep 2024)-img/p1_phy_s1_q04_1.png}{62.5}{0}
\qfigdef{p1_phy_s1_q05_1}{SRG (Enthusiast) - Test 39 (15 Sep 2024)-img/p1_phy_s1_q05_1.png}{99.9}{0}
\qfigdef{p1_phy_s1_q06_1}{SRG (Enthusiast) - Test 39 (15 Sep 2024)-img/p1_phy_s1_q06_1.png}{105.8}{0}
\qfigdef{p1_phy_s2_q04_1}{SRG (Enthusiast) - Test 39 (15 Sep 2024)-img/p1_phy_s2_q04_1.png}{118.5}{0}
\qfigdef{p1_phy_s2_q06_1}{SRG (Enthusiast) - Test 39 (15 Sep 2024)-img/p1_phy_s2_q06_1.png}{107.3}{0}
\qfigdef{p1_phy_s3_q01_1}{SRG (Enthusiast) - Test 39 (15 Sep 2024)-img/p1_phy_s3_q01_1.png}{87.2}{0}
\qfigdef{p1_phy_s3_q03_1}{SRG (Enthusiast) - Test 39 (15 Sep 2024)-img/p1_phy_s3_q03_1.png}{142.3}{0}
\qfigdef{p1_phy_s3_q04_1}{SRG (Enthusiast) - Test 39 (15 Sep 2024)-img/p1_phy_s3_q04_1.png}{108.3}{0}
\qfigdef{p1_phy_s3_q05_1}{SRG (Enthusiast) - Test 39 (15 Sep 2024)-img/p1_phy_s3_q05_1.png}{303.9}{1}
\qfigdef{p1_phy_s3_q06_1}{SRG (Enthusiast) - Test 39 (15 Sep 2024)-img/p1_phy_s3_q06_1.png}{158.8}{0}
\qfigdef{p1_chem_s1_q01_1}{SRG (Enthusiast) - Test 39 (15 Sep 2024)-img/p1_chem_s1_q01_1.png}{128.3}{0}
\qfigdef{p1_chem_s1_q02_1}{SRG (Enthusiast) - Test 39 (15 Sep 2024)-img/p1_chem_s1_q02_1.png}{185.2}{0}
\qfigdef{p1_chem_s1_q06_A}{SRG (Enthusiast) - Test 39 (15 Sep 2024)-img/p1_chem_s1_q06_A.png}{259.6}{1}
\qfigdef{p1_chem_s1_q06_2}{SRG (Enthusiast) - Test 39 (15 Sep 2024)-img/p1_chem_s1_q06_2.png}{63.4}{0}
\qfigdef{p1_chem_s1_q06_C}{SRG (Enthusiast) - Test 39 (15 Sep 2024)-img/p1_chem_s1_q06_C.png}{237.1}{0}
\qfigdef{p1_chem_s2_q02_1}{SRG (Enthusiast) - Test 39 (15 Sep 2024)-img/p1_chem_s2_q02_1.png}{176.8}{0}
\qfigdef{p1_chem_s2_q03_1}{SRG (Enthusiast) - Test 39 (15 Sep 2024)-img/p1_chem_s2_q03_1.png}{230.2}{0}
\qfigdef{p1_chem_s2_q04_1}{SRG (Enthusiast) - Test 39 (15 Sep 2024)-img/p1_chem_s2_q04_1.png}{107.3}{0}
\qfigdef{p1_chem_s3_q01_1}{SRG (Enthusiast) - Test 39 (15 Sep 2024)-img/p1_chem_s3_q01_1.png}{159.9}{0}
\qfigdef{p1_chem_s3_q03_1}{SRG (Enthusiast) - Test 39 (15 Sep 2024)-img/p1_chem_s3_q03_1.png}{116.8}{0}
\qfigdef{p1_chem_s3_q05_1}{SRG (Enthusiast) - Test 39 (15 Sep 2024)-img/p1_chem_s3_q05_1.png}{131.8}{0}
\qfigdef{p1_chem_s3_q06_1}{SRG (Enthusiast) - Test 39 (15 Sep 2024)-img/p1_chem_s3_q06_1.png}{225.5}{0}

\begin{document}
\begin{center}{\Large\bfseries SRG (Enthusiast) --- Test 39 (15 Sep 2024)}\end{center}
\begin{multicols}{2}
\partsec{1}{PHYSICS}
\sect{Section-I}
\qn{1} The position vector of a particle that is moving in space is given by\par
$\vec r=(1+2\cos2\omega t)\,\hat i+(3\sin^2\omega t)\,\hat j+(3)\,\hat k$\par
in the ground frame. All units are in SI. Choose the correct statement (s) :
\begin{qoptions}
\opt{A}{The particle executes SHM in the ground frame about the mean position $\left(1,\dfrac32,3\right)$}
\opt{B}{The particle executes SHM the ground frame about mean position $\left(1,\dfrac32,0\right)$}
\opt{C}{The amplitude of the SHM of the particle is $\dfrac52$ m.}
\opt{D}{Particle does not perform SHM}
\end{qoptions}
\qn{2} Two objects P and Q of masses $2m$ and $m$ are connected to two springs of spring constant $2k$ and $k$ respectively. Other ends of springs are connected to rigid walls as shown. Initially springs are compressed by lengths $2A$ and $A$ respectively, and then released. (before collision time period is calculated in absence of other block)
\qfig{p1_phy_s1_q02_1}
\begin{qoptions}
\opt{A}{Time period of P and Q remains same before and after the collision irrespective of the fact whether collision between masses is perfectly elastic or perfectly inelastic.}
\opt{B}{Time period of P and Q remains same before and after the collision only if collision between the objects is perfectly elastic.}
\opt{C}{Time period of P and Q remains same before and after the collision only if collision between the objects is perfectly inelastic.}
\opt{D}{If collision is perfectly inelastic, amplitude of S.H.M. of combined mass is $A$.}
\end{qoptions}
\qn{3} In the figure below, the switches $S_1$ and $S_2$ are closed simultaneously at $t=0$ and a current starts to flow in the circuit. Both the batteries have the same magnitude of the electromotive force (\textit{emf}) and the polarities are as indicated in the figure. Ignore mutual inductance between the inductors. The current $I$ in the middle wire reaches its maximum magnitude $I_{max}$ at time $t=\tau$. Which of the following statement(s) is (are) true?
\qfig{p1_phy_s1_q03_1}
\begin{qoptions}
\opt{A}{$I_{\max}=\dfrac{V}{2R}$}
\opt{B}{$I_{\max}=\dfrac{V}{4R}$}
\opt{C}{$\tau=\dfrac{L}{R}\ell\mathrm n2$}
\opt{D}{$\tau=\dfrac{2L}{R}\ell\mathrm n2$}
\end{qoptions}
\qn{4} In the cylindrical region of radius $R=2$\,m, there exists a time varying magnetic field $B$ such that $\dfrac{dB}{dt}=2$ Tesla/sec. A charged particle having charge $q=2$\,C is placed at the point P at a distance $d$ from its centre O. Now, the particle is moved in the direction perpendicular to line OP by an external agent to infinity so that there is no gain in kinetic energy of charged particle. Then
\qfig{p1_phy_s1_q04_1}
\begin{qoptions}
\opt{A}{Work done by external agent is $4\pi$ joule when $d=4$\,m}
\opt{B}{Work done by external agent is $4\pi$ Joule when $d=8$ m}
\opt{C}{Work done by external agent is independent of $d$ (as long as $d$ is more than $R$)}
\opt{D}{Work done by external agent is positive}
\end{qoptions}
\qn{5} A conducting wire of parabolic shape, initially $y=x^2$, is moving with velocity $\vec V=V_0\hat i$ in a magnetic field $\vec B=B_0\left[1+\left(\dfrac{y}{L}\right)^{\beta}\right]\hat k$, $B_0$, $L$ and $\beta$ are positive constants and $\Delta\phi$ is the potential difference developed between the ends of the wire, then the correct statement(s) is/are:
\qfig{p1_phy_s1_q05_1}
\begin{qoptions}
\opt{A}{$|\Delta\phi|$ remains the same if the parabolic wire is replaced by a straight wire, $y=x$ initially, of length $\sqrt2L$.}
\opt{B}{$|\Delta\phi|$ is proportional to the length of the wire projected on the $y$-axis.}
\opt{C}{$|\Delta\phi|=\dfrac12B_0V_0L$ for $\beta=0$.}
\opt{D}{$|\Delta\phi|=\dfrac43B_0V_0L$ for $\beta=2$.}
\end{qoptions}
\qn{6} Two parallel resistanceless rails are connected by an inductor of inductance $L$ at one end as shown in the figure. A magnetic field $B$ exists in the space which is perpendicular to the plane of the rails. Now a conductor of length $\ell$ and mass $m$ is placed transverse on the rails and given an impulse $J$ towards the rightward direction. Then choose the \textbf{correct} option (s). [ignore resistance of conductor]
\qfig{p1_phy_s1_q06_1}
\begin{qoptions}
\opt{A}{Velocity of the conductor is half of the initial velocity after a displacement of the conductor $d=\sqrt{\dfrac{3J^2L}{4B^2\ell^2m}}$}
\opt{B}{Current flowing through the inductor at the instant when velocity of the conductor is half of the initial velocity is $i=\sqrt{\dfrac{3J^2}{4Lm}}$}
\opt{C}{Velocity of the conductor is half of the initial velocity after a displacement of the conductor $d=\sqrt{\dfrac{3J^2L}{B^2\ell^2m}}$}
\opt{D}{Current flowing through the inductor at the instant when velocity of the conductor is half of the initial velocity is $i=\sqrt{\dfrac{3J^2}{mL}}$}
\end{qoptions}
\sect{Section-II}
\qn{1} Quality factor of a series LCR circuit is 0.6. If definition of quality factor is changed such that now limiting frequency of bandwidth is taken upto $\dfrac13$rd of maximum power rather than half of maximum power (standard definition), then what will be quality factor of given circuit ?
\qn{2} In R-L-C series circuit, we have same current at angular frequencies $\omega_1$ and $\omega_2$. Find the resonant frequency.\par
(Take $\omega_1=4$ units, $\omega_2=9$ units)
\qn{3} A long coaxial cable consists of an inner conductor of radius `$a$' and outer conductor of radius `$3a$'. The two conductors carry equal current $I$ in opposite directions. Magnetic energy stored per unit length of the cable is $\dfrac{\mu_0I^2}{4\pi}\ell\mathrm n(x)$. The value of $x$ is :
\qn{4} The circuit consists of a capacitor, inductance, and two resistors, see figure. The voltage on both resistors is 10\,V, and the voltage between the points A and B is also 10\,V. Find rms value of the applied voltage $U_0$ (in volt).
\qfig{p1_phy_s2_q04_1}
\qn{5} An inductor ($L$), a resistance ($R$) and a ideal battery (emf $\varepsilon$) are connected in series. Circuit is switched on at $t=0$. The ratio of energy stored in the inductor ($U_1$) at time $t=L/R$ and energy dissipated in resistance ($U_2$) up to that instant is $\dfrac{U_1}{U_2}=\dfrac12\,\dfrac{\left(1-\frac1e\right)^2}{\left(\frac{\beta}{e}-\frac{1}{\gamma e^2}-\frac1\delta\right)}$. Here $\alpha$, $\beta$, $\gamma$ and $\delta$ are integers, $e$ is Euler's number. Find $(\beta+\gamma+\delta)$.
\qn{6} A block is attached with spring (1) as shown in figure. An identical spring is also attached with fixed wall on other side and it's free end is at a distance $(A/2)$ from block. Now block is compressed by distance $A$ against spring (1) and released. If the time period of oscillation of block is $T=\sqrt{\dfrac mk}\left[\dfrac{x\pi}{y}+\sqrt2\cos^{-1}\left(\dfrac{1}{\sqrt z}\right)\right]$ (where, $x$, $y$ are co-prime). Find $x+y+z$.
\qfig{p1_phy_s2_q06_1}
\sect{Section-III}
\qn{1} A loop is formed by two parallel conductors connected by a solenoid with inductance $L=2$\,H and a conducting rod of mass $m=80$\,g, which can move freely (without friction) slide over the conductors. The conductors are located in a horizontal plane with a uniform vertical magnetic field $B=4$\,T. The distance between conductors is equal to $\ell=4$\,cm. At the moment $t=0$, the rod is imparted an initial velocity $v_0=2$ cm/s directed to the right. The amplitude (in cm) for the oscillations of conducting rod.
\qfig{p1_phy_s3_q01_1}
\qn{2} In a series L-C-R circuit $R=12\,\Omega$, $X_L$(at resonance) $=6\,\Omega$ is operated at an frequency which is half of resonance frequency $\left(\omega=\dfrac{\omega_R}{2}\right)$ then the value of power factor of the circuit is $\dfrac X5$ find $X$
\qn{3} In the circuit shown, switch $S_2$ is open and $S_1$ is closed since long. Take $E=20$\,V, $L=0.5$ H and $R=10\,\Omega$. The magnitude of rate of change of energy stored in the magnetic field inside the inductor, immediately after $S_2$ is closed is `$N$' J/s then $\dfrac{N}{10}$ is \underline{\hspace{1.5em}}.
\qfig{p1_phy_s3_q03_1}
\qn{4} A resistance less coil is connected in parallel to a resistor $R$ through switch K to a source with constant EMF $E$ and internal resistance $r$. At the initial moment of time the key is open and there is no current in the circuit. What charge $Q$ (in mC) will flow through resistance $R$ after closing key K? Neglect the ohmic resistance of the coil. Take $E=6$ volts, $L=3$\,mH, $R=2\,\Omega$, $r=3\,\Omega$.
\qfig{p1_phy_s3_q04_1}
\qfullbegin
\qn{5} If this body rolls purely then time period of its small oscillations is $T$. It's moment of inertia about CM is 9\,kg.m$^2$.\par
Find value of $\dfrac{40}{\pi}T$ ?
\qfig{p1_phy_s3_q05_1}
\qfullend
\qn{6} Two blocks 1 and 2, each having mass $m$, are placed on a smooth table connected to three identical springs as shown in the figure. Each spring has a force constant $k$. Initially, all springs are relaxed. The system is disturbed and starts moving. Let $x_1$ and $x_2$ represent the displacements of the two blocks from their respective mean positions. Motion of block 1 is superposition of two SHMs. If the angular frequency of the component SHMs is $\omega_a$ \& $\omega_b$ and if value of $\omega_a\times\omega_b=\sqrt\alpha\,\dfrac km$, fill $\alpha$.
\qfig{p1_phy_s3_q06_1}
\partsec{2}{CHEMISTRY}
\sect{Section-I}
\qn{1} The reaction between NO$_2$ and F$_2$ gas:\par
2NO$_2$(g) + F$_2$(g) $\to$ 2NO$_2^+$(g); Rate $=K[\mathrm{NO_2}][\mathrm{F_2}]$\par
The accepted mechanism for the reaction is\par
Step $-$ 1: NO$_2$(g) + F$_2$(g) $\to$ NO$_2$F(g) + F(g)\par
Step $-$ 2: NO$_2$(g) + F(g) $\to$ NO$_2$F(g)\par
Which of the following statements are true?
\qfig{p1_chem_s1_q01_1}
\begin{qoptions}
\opt{A}{Order of reaction $=$ molecularity at R.D.S}
\opt{B}{Reaction is exothermic}
\opt{C}{Step $-$ 1 is rate determing step}
\opt{D}{Number of transition state $=1$}
\end{qoptions}
\qn{2} The following compounds are basic in nature (H$^+-$acceptor)
\qfig{p1_chem_s1_q02_1}
The correct decreasing order of basicity of the given amines is:
\begin{qoptions}
\opt{A}{5 $>$ 3 $>$ 6 $>$ 4}
\opt{B}{2 $>$ 1 $>$ 5 $>$ 3}
\opt{C}{1 $>$ 5 $>$ 3 $>$ 4}
\opt{D}{6 $>$ 2 $>$ 5 $>$ 1}
\end{qoptions}
\qn{3} Which of the following statement(s) are \textbf{INCORRECT} about compounds of ``Boron family'' elements ?
\begin{qoptions}
\opt{A}{Diborane, react spontaneously with air, often with explosive violence and green flash, the final product is boric anhydride}
\opt{B}{BCl$_3$, BBr$_3$ and BI$_3$ are susceptible to hydrolysis}
\opt{C}{Lewis acidity : AlCl$_3>$ BCl$_3>$ GaCl$_3$}
\opt{D}{In borazine compound, the $\pi$ electrons are concentrated on `B' atoms and there is a net partial positive charge on N-atoms}
\end{qoptions}
\qn{4} Which of the following statements are \textbf{CORRECT} about compounds of CARBON FAMILY elements?
\begin{qoptions}
\opt{A}{Electrical conductivity of graphite increases with temperature.}
\opt{B}{Silane (SiH$_4$) is prepared by the reduction of SiO$_2$ with Al under high pressure of H$_2$.}
\opt{C}{`CO' gas is used in the production of H$_2$ through water-gas shift reaction.}
\opt{D}{Hydrolysis of CH$_3$SiCl$_3$ followed by polymerisation produces a cross-linked polymers which contains ``Si-O-Si'' bonds}
\end{qoptions}
\qn{5} Which of the following statements are \textbf{CORRECT} about ``Radioactive'' elements?
\begin{qoptions}
\opt{A}{If the nuclide is too heavy to be stable [$Z>83$], it may undergoes `$\alpha$' (alpha) decay.}
\opt{B}{Due to electron (e$^-$) capture, element position shifts in left direction in periodic table.}
\opt{C}{Radioactive decay follows first order kinetics}
\opt{D}{If $A$ = Activity of radio active substance then : $A_t=A_oe^{-kt}$; $k$ = decay constant}
\end{qoptions}
\qfullbegin
\qn{6} Which of the following statements(s) is/are \textbf{CORRECT} about amines ?
\begin{qoptions}
\opt{A}{\qfigi{p1_chem_s1_q06_A}}
\opt{B}{Boiling point: \qfigi{p1_chem_s1_q06_2} $<$ (Et)$_2$NH $<$ CH$_3-$CH$_2-$N(CH$_3$)$_2$ (in aq. solution)}
\opt{C}{\qfigi{p1_chem_s1_q06_C}}
\opt{D}{Bromination of aniline in H$_2$O gives 2, 4 - dibromoaniline as major product}
\end{qoptions}
\qfullend
\sect{Section-II}
\qn{1} Number of B $-$ O $-$ B bonds in borax is $\mathbf x$ then value of $\dfrac x5$ is?
\qn{2} In the following reaction
\qfig{p1_chem_s2_q02_1}
The percentage yield of the product (\textbf{P}) is 100\%. The weight in grams of the product (\textbf{P}) formed by completely reacting 21.25 grams of reactant (\textbf{R}) is
\qn{3} Number of bromine atoms in one molecule of the product (\textbf{P}) is
\qfig{p1_chem_s2_q03_1}
\qn{4} The azo-dye (\textbf{Y}) formed in the following reactions is\par
Sulphanilic acid $+$ NaNO$_2$ $+$ CH$_3$COOH $\to$ X
\qfig{p1_chem_s2_q04_1}
If number of sp$^2-$hybridized CARBON atoms $=$ `a' (in Y)\par
number of N$-$atoms $=$ `b' (in Y)\par
then value $\left(\dfrac ab\right)=$ \underline{\hspace{3em}}?
\qn{5} For the reaction A $\to$ B, the rate constant $k$ (in s$^{-1}$) is given by $\log_{10}k=20.35-\dfrac{\left(2.47\times10^{3}\right)}{T}$. The energy of activation in kJ mole$^{-1}$ is \underline{\hspace{2em}}.\par
[Given: $R=8.314$ JK$^{-1}$ mol$^{-1}$]
\qn{6} The total number of `$\alpha$' and `$\beta$' particles emitted in the nuclear reaction $^{238}_{\ 92}$U $\to$ $^{214}_{\ 82}$Pb is \underline{\hspace{3em}} ?
\sect{Section-III}
\qn{1} Consider the following representation of the reaction\par
2AB$_2\to$ 2AB + B$_2$
\qfig{p1_chem_s3_q01_1}
If the given reaction follows zero order kinetics, the value of $\dfrac x5=$ \underline{\hspace{3em}}?
\qn{2} Consider the general reaction: aA + bB $\to$ cC and the following average rate data over some time period $\Delta t$ :\par
$-\dfrac{\Delta A}{\Delta t}=0.0080$ mol/L.\,s\par
$-\dfrac{\Delta B}{\Delta t}=0.0120$ mol/L.\,s\par
$\dfrac{\Delta C}{\Delta t}=0.0160$ mol/L.\,s\par
If the coefficient of A is 2, then sum of $(a+b+c)=$ \underline{\hspace{3em}}
\qn{3} According to following reaction scheme:
\qfig{p1_chem_s3_q03_1}
(A and D) = Boron containing products.\par
(E and C) = Same gaseous product.\par
If $x$ = maximum no. of moles of CH$_3$OH required for number of moles of `A' formed.\par
$\mathbf y$ = total number of mole of A, C, D and E, formed in the above reaction sequence then value of $|\mathbf{y-x}|$ is \underline{\hspace{3em}} ?
\qn{4} The primary unit of cyclic metasilicate is $\left[(\mathrm{SiO}_3)_n\right]^{-2n}$ and mineral benitoite is an example:\par
Benitoite: Ba$_x$Ti$_y$Si$_3$O$_z$\par
The sum of value of $(z-y-x)$ is \underline{\hspace{3em}} ?
\qn{5} \qfig{p1_chem_s3_q05_1}
Number of ``PRIMARY CARBON'' atoms in product `\textbf{S}' \underline{\hspace{3em}} ?
\qn{6} According to given reaction scheme:
\qfig{p1_chem_s3_q06_1}
Degree of unsaturation of product (P) is \underline{\hspace{3em}} ?
\partsec{3}{MATHEMATICS}
\sect{Section-I}
\qn{1} Consider the function $f(k)=\displaystyle\int_0^{\pi/2}\dfrac{dx}{\sqrt{1-k\cos^2x}}$ where $k\in(0,1)$, then which of the following is/are correct-
\begin{qoptions}
\opt{A}{$f\!\left(\frac12\right)<f\!\left(\frac13\right)<f\!\left(\frac14\right)$}
\opt{B}{$f\!\left(\frac14\right)<f\!\left(\frac13\right)<f\!\left(\frac12\right)$}
\opt{C}{$f\!\left(\frac14\right)<f\!\left(\frac12\right)>f\!\left(\frac34\right)$}
\opt{D}{$f\!\left(\frac14\right)<f\!\left(\frac12\right)<f\!\left(\frac34\right)$}
\end{qoptions}
\qn{2} Let $A_n=\displaystyle\sum_{k=0}^{n-1}\dfrac{n}{k^2+n^2+(n+2)k+n+1}$ and $B_n=\displaystyle\sum_{k=0}^{n-1}\dfrac{n}{n^2+kn+k^2}$ for $n=1,2,3,\ldots$, then which of the following is/are true ?
\begin{qoptions}
\opt{A}{$A_n<B_n$}
\opt{B}{$A_n>B_n$}
\opt{C}{$A_n<\dfrac\pi3$}
\opt{D}{$B_n>\dfrac\pi9$}
\end{qoptions}
\qn{3} Consider the equation $x\,|x|+px+q=0$, $p,q\in R$. Then which of the following is/are correct
\begin{qoptions}
\opt{A}{Equation has atmost 3 real solutions}
\opt{B}{Equation has atleast 1 real solutions}
\opt{C}{Equation has real solutions only if $p^2-4q\ge0$}
\opt{D}{Equation has 3 real solutions if $p<0$ and $q>0$.}
\end{qoptions}
\qn{4} Let $\alpha$ and $\beta$ be the roots $x^2-x-1=0$, with $\alpha>\beta$. For all positive integers $n$, define\par
$a_n=\dfrac{\alpha^n-\beta^n}{\alpha-\beta}$, $n\ge1$\par
$b_1=1$ and $b_n=a_{n-1}+a_{n+1}$, $n\ge2$.\par
Then which of the following options is/are correct?
\begin{qoptions}
\opt{A}{$a_1+a_2+a_3+\ldots+a_n=a_{n+2}-1$ for all $n\ge1$.}
\opt{B}{$\displaystyle\sum_{n=1}^{\infty}\dfrac{a_n}{10^n}=\dfrac{10}{89}$}
\opt{C}{$\displaystyle\sum_{n=1}^{\infty}\dfrac{b_n}{10^n}=\dfrac{12}{89}$}
\opt{D}{$b_n=\alpha^n+\beta^n$ for all $n\ge1$.}
\end{qoptions}
\qn{5} Consider the equation $(x^2-4x+6)^2-(k+2)(x^2-4x+6)+k+5=0$ ($k\in R$). Which of following are correct
\begin{qoptions}
\opt{A}{There is no integral value of $k$ for which given equation has 4 distinct real solutions.}
\opt{B}{There is only one integral value of $k$ for which given equation has 3 distinct real solutions.}
\opt{C}{There is only one integral value of $k$ for which given equation has exactly one real solution.}
\opt{D}{The number of positive integral values of $k$ for which given equation has no real solution is 3.}
\end{qoptions}
\qn{6} Let $f:(0,\infty)\to R$ be given by $f(x)=\displaystyle\int_{1/x}^{x}e^{-\left(t+\frac1t\right)}\dfrac{dt}{t}$. Then
\begin{qoptions}
\opt{A}{$f(x)$ is monotonically increasing on $[1,\infty)$}
\opt{B}{$f(x)$ is monotonically decreasing on $(0,1)$}
\opt{C}{$f(x)+f\!\left(\dfrac1x\right)=0$, for all $x\in(0,\infty)$}
\opt{D}{$f(2^x)$ is an odd function of $x$ on $R$}
\end{qoptions}
\sect{Section-II}
\qn{1} If $\alpha,\beta,\gamma$ are roots of $x^3-2x^2+3x-2=0$, then the value of $\left(\dfrac{\alpha\beta}{\alpha+\beta}+\dfrac{\alpha\gamma}{\alpha+\gamma}+\dfrac{\beta\gamma}{\beta+\gamma}\right)$ is
\qn{2} Given $f(x)=x^2+\alpha x+\beta$ (where $\alpha,\beta\in R$). If the equation $f(f(x))=0$ has two roots 1 and 2 such that $f(1)\neq f(2)$, then $|f(0)|$ is equal to $-$
\qn{3} Let $\alpha,\beta$ ($\alpha>\beta$) be the roots of the quadratic equation $x^2-x-2=0$. If $P_n=\alpha^n-\beta^n$, $n\in N$, then $\dfrac{(P_{10}-P_9)(P_9-P_8)}{P_7\cdot P_8}$ is equal to
\qn{4} If $\displaystyle\lim_{n\to\infty}\left(\dfrac{{}^{3n}C_n}{{}^{2n}C_n}\right)^{1/n}=\dfrac ab$ where $a,b$ are relatively prime positive integers then $a+b=$
\qn{5} Let $f(x)=\max\{|x+1|,|x+2|,\ldots,|x+5|\}$. Then $\displaystyle\int_{-6}^{0}f(x)\,dx$ is equal to \underline{\hspace{4em}}.
\qn{6} The number of values of $x$ satisfying $\displaystyle\int_0^{x}t^2\sin(x-t)\,dt=x^2$ in $[0,100]$ is
\sect{Section-III}
\qn{1} Let $P(x)=x^2-x+1$ and $\alpha$ be a root of the equation $P(P(P(P(x))))=0$, then the absolute value of $(P(\alpha)-1)\,P(\alpha)\,P(P(\alpha))\,P(P(P(\alpha)))$ is\ldots\ldots\ldots
\qn{2} The number of values of $K$ for which the expression $\left(2x^2-(K+1)x+4\right)\times\left(x^2-Kx+(K+1)\right)$ is a perfect square of quadratic polynomial is \underline{\hspace{1.5em}}
\qn{3} Let $f(x)=x^2+ax+b$. If the maximum and the minimum values of $f(x)$ are 3 and 2 respectively for $0\le x\le2$, then the number of possible ordered pair(s) of $(a,b)$ is :
\qn{4} If $\displaystyle\sum_{n=1}^{\infty}\left(\dfrac1{4n-3}-\dfrac1{4n-1}\right)=\dfrac\pi k$, then $k=$ \underline{\hspace{2.5em}}
\qn{5} If $f(\pi)=2$ \& $\displaystyle\int_0^{\pi}\left(f(x)+f''(x)\right)\sin x\,dx=5$, then $f(0)=$ \underline{\hspace{2.5em}}
\qn{6} The value of $\displaystyle\int_0^{1}\sqrt[3]{2x^3-3x^2-x+1}\,dx$ is $-$
\end{multicols}
\end{document}
